% GENERATED FILE. Edit canonical lessons, then run npm run generate:books.
% canonical-source-hash: fnv1a64:284897490760e160
% canonical-lessons: GE-C195-fussgaengerzone, GE-C195-dom, GE-C195-denkmal, GE-C195-turm, GE-C195-burg

\chapter{A pedestrian zone, A cathedral, A monument, A tower, A castle}
\label{ch:ge-a-pedestrian-zone-a-cathedral-a-monument-a-tower-a-castle}
\hlchaptermodality{195}

\begin{chapteropening}
\textbf{By the end of this chapter:} \emph{I can say a pedestrian zone, a cathedral, a monument, a tower and a castle.}

Complete the last lesson of chapter 195: I can say a pedestrian zone, a cathedral, a monument, a tower and a castle.
\end{chapteropening}

\section[die Fußgängerzone]{die Fußgängerzone — a pedestrian zone}

\label{lesson:GE-C195-fussgaengerzone}



\begin{quote}
Before the new one: say the German for a round trip, then the German for a market square.
\end{quote}

\subsection*{You'll want to know: die Fußgängerzone}
\textbf{die Fußgängerzone} — \textquotedblleft{}a pedestrian zone\textquotedblright{}.

\begin{cousinweb}
\textbf{Fußgängerzone} is \emph{Fuß} + a noun built on \emph{gehen}, \textquotedblleft{}to go\textquotedblright{}, + \emph{Zone}, \textquotedblleft{}zone\textquotedblright{}, which came through Latin from Greek \emph{z\'{\={o}}nē}, \textquotedblleft{}belt\textquotedblright{}: a belt of town for people on foot.
\end{cousinweb}

\begin{grammarlens}[title={What you've built}]
One more word for this chapter.
\end{grammarlens}

\subsection*{Your turn}
\begin{itemize}
\raggedright
  \item \emph{Say it:} \emph{die Fußgängerzone}
  \item \emph{Say it:} \emph{die Fußgängerzone}, once more
  \item \emph{Say it:} say \emph{die Fußgängerzone}
  \item \emph{Recall:} say the German for a round trip, then the German for a market square, then say \emph{die Fußgängerzone} again
\end{itemize}

\begin{tcolorbox}[breakable,colback=teal!4,colframe=teal!35!black,title={Before you move on}]
What does die Fußgängerzone mean? (\textquotedblleft{}A pedestrian zone\textquotedblright{}.) Say it once more.
\end{tcolorbox}
\section[der Dom]{der Dom — a cathedral}

\label{lesson:GE-C195-dom}



\begin{quote}
Before the new one: say the German for a market square, then the German for a pedestrian zone.
\end{quote}

\subsection*{You'll want to know: der Dom}
\textbf{der Dom} — \textquotedblleft{}a cathedral\textquotedblright{}.

\begin{cousinweb}
\textbf{Dom} $\leftarrow$ Middle High German \emph{tuom} $\leftarrow$ Medieval Latin \emph{domus} (\emph{episcopalis}), \textquotedblleft{}the (bishop's) house\textquotedblright{} — borrowed. The modern form was reshaped after French \emph{dôme} (EtymDB 2.1).
\end{cousinweb}

\begin{grammarlens}[title={What you've built}]
One more word for this chapter.
\end{grammarlens}

\subsection*{Your turn}
\begin{itemize}
\raggedright
  \item \emph{Say it:} \emph{der Dom}
  \item \emph{Say it:} \emph{der Dom}, once more
  \item \emph{Say it:} say \emph{der Dom}
  \item \emph{Recall:} say the German for a market square, then the German for a pedestrian zone, then say \emph{der Dom} again
\end{itemize}

\begin{tcolorbox}[breakable,colback=teal!4,colframe=teal!35!black,title={Before you move on}]
What does der Dom mean? (\textquotedblleft{}A cathedral\textquotedblright{}.) Say it once more.
\end{tcolorbox}
\section[das Denkmal]{das Denkmal — a monument}

\label{lesson:GE-C195-denkmal}



\begin{quote}
Before the new one: say the German for a pedestrian zone, then the German for a cathedral.
\end{quote}

\subsection*{You'll want to know: das Denkmal}
\textbf{das Denkmal} — \textquotedblleft{}a monument\textquotedblright{}.

\begin{cousinweb}
\textbf{Denkmal} is \emph{denken}, \textquotedblleft{}to think, to remember\textquotedblright{}, + the old word for \textquotedblleft{}a mark\textquotedblright{}, Old High German \emph{māl} (EtymDB 2.1): a mark to remember by, copying Greek \emph{mnēmósynon}.
\end{cousinweb}

\begin{grammarlens}[title={What you've built}]
One more word for this chapter.
\end{grammarlens}

\subsection*{Your turn}
\begin{itemize}
\raggedright
  \item \emph{Say it:} \emph{das Denkmal}
  \item \emph{Say it:} \emph{das Denkmal}, once more
  \item \emph{Say it:} say \emph{das Denkmal}
  \item \emph{Recall:} say the German for a pedestrian zone, then the German for a cathedral, then say \emph{das Denkmal} again
\end{itemize}

\begin{tcolorbox}[breakable,colback=teal!4,colframe=teal!35!black,title={Before you move on}]
What does das Denkmal mean? (\textquotedblleft{}A monument\textquotedblright{}.) Say it once more.
\end{tcolorbox}
\section[der Turm]{der Turm — a tower}

\label{lesson:GE-C195-turm}



\begin{quote}
Before the new one: say the German for a cathedral, then the German for a monument.
\end{quote}

\subsection*{You'll want to know: der Turm}
\textbf{der Turm} — \textquotedblleft{}a tower\textquotedblright{}.

\begin{cousinweb}
\textbf{Turm} $\leftarrow$ Middle High German \emph{turn} $\leftarrow$ Old French \emph{torn}, \emph{tor} $\leftarrow$ Latin \emph{turris}, \textquotedblleft{}tower\textquotedblright{} (EtymDB 2.1) — borrowed, like English \emph{tower}.
\end{cousinweb}

\begin{grammarlens}[title={What you've built}]
One more word for this chapter.
\end{grammarlens}

\subsection*{Your turn}
\begin{itemize}
\raggedright
  \item \emph{Say it:} \emph{der Turm}
  \item \emph{Say it:} \emph{der Turm}, once more
  \item \emph{Say it:} say \emph{der Turm}
  \item \emph{Recall:} say the German for a cathedral, then the German for a monument, then say \emph{der Turm} again
\end{itemize}

\begin{tcolorbox}[breakable,colback=teal!4,colframe=teal!35!black,title={Before you move on}]
What does der Turm mean? (\textquotedblleft{}A tower\textquotedblright{}.) Say it once more.
\end{tcolorbox}
\section[die Burg]{die Burg — a castle}

\label{lesson:GE-C195-burg}



\begin{quote}
Before the new one: say the German for a monument, then the German for a tower.
\end{quote}

\subsection*{You'll want to know: die Burg}
\textbf{die Burg} — \textquotedblleft{}a castle\textquotedblright{}.

\begin{cousinweb}
\textbf{Burg} $\leftarrow$ Old High German \emph{burg} $\leftarrow$ Proto-Germanic \emph{*burgz}, \textquotedblleft{}stronghold\textquotedblright{} — inherited. (EtymDB 2.1) English \emph{borough} and the \emph{-bury} of place names are the same word.
\end{cousinweb}

\begin{grammarlens}[title={What you've built}]
One more word for this chapter.
\end{grammarlens}

\subsection*{Your turn}
\begin{itemize}
\raggedright
  \item \emph{Say it:} \emph{die Burg}
  \item \emph{Say it:} \emph{die Burg}, once more
  \item \emph{Say it:} say \emph{die Burg}
  \item \emph{Recall:} say the German for a monument, then the German for a tower, then say \emph{die Burg} again
\end{itemize}

\begin{tcolorbox}[breakable,colback=teal!4,colframe=teal!35!black,title={Before you move on}]
What does die Burg mean? (\textquotedblleft{}A castle\textquotedblright{}.) Say it once more.
\end{tcolorbox}
